\documentclass[12pt]{report}
\usepackage[utf8]{inputenc}
\usepackage[T1]{fontenc}
\usepackage{geometry}
\usepackage{hyperref}
\usepackage{amsmath}
\geometry{margin=1in}
\title{Atelier Long Report Fixture}
\author{Atelier Test Suite}
\date{\today}

\begin{document}
\maketitle
\tableofcontents

\chapter{Project Overview}
This report is a second long document for simultaneous-tab testing. It uses the
report class so that compilation produces a different outline hierarchy and a
larger PDF with chapter boundaries.

\section{Scope}
The soak covers source editing, preview resizing, compilation, diagnostics,
comments, outline navigation, history, and synchronization. It intentionally
does not depend on external bibliography or image files.

\section{Expected workflow}
Open this report beside the long article. Keep the preview visible in one tab
and hidden in the other. Resize both, reload them, and confirm that the stored
split preference behaves consistently.

\chapter{Long Narrative}
\section{First sequence}
The first sequence supplies several paragraphs for scrolling and wrapping. Each
paragraph is long enough to span multiple visual lines at ordinary editor widths.

During editing, the line number gutter should remain aligned with the source.
Changing the split width must trigger a CodeMirror refresh without moving the
cursor or resetting the scroll position.

The preview panel should remain usable when the browser window becomes narrow.
Neither pane should collapse completely, and the separator should continue to
indicate that horizontal dragging is available.

\section{Second sequence}
Background tasks may finish in a different order from the actions that started
them. The visible status should prioritize save, compile, conflict, and modified
states over secondary persistence notifications.

The unique phrase REPORT-COMMENT-ANCHOR can be selected for a persistent
annotation. Rewrap this paragraph and reload the file to validate re-anchoring.

\section{Third sequence}
Introduce an invalid command on a new line, compile, inspect the gutter marker,
and click the corresponding log entry. Remove the invalid command and compile
again; the diagnostic decoration should disappear.

\chapter{Repeated Editing Material}
\section{Iteration one}
Iteration one represents an early drafting pass. Add a sentence, save it, and
inspect the version history. The editor should report a saved state only after
the server confirms the write.

\section{Iteration two}
Iteration two represents a revision pass. Select several lines, replace them,
undo the replacement, and redo it. The history should remain local to this tab.

\section{Iteration three}
Iteration three represents a formatting pass. Change the wrap column, rewrap a
comment block, and ensure that ordinary LaTeX commands remain untouched.

\section{Iteration four}
Iteration four represents a review pass. Add an anchored comment, navigate away,
return through the outline, and verify that the highlighted range still matches
the selected source phrase.

\section{Iteration five}
Iteration five represents a compilation pass. Modify the source without saving
and press Compile. The shell should persist the buffer before invoking LaTeX.

\section{Iteration six}
Iteration six represents a synchronization pass. Move from source to PDF, then
back from the PDF coordinates to the source line. The cursor highlight should be
brief and should not become a permanent decoration.

\chapter{Technical Notes}
\section{Split state}
The split percentage is stored under a dedicated editor preference. Values are
bounded so that the editor and preview both retain a useful minimum width.

\section{Lifecycle}
Every listener added by the LaTeX module must be removed when the tab closes.
This includes cursor events, input events, message events, separator drag events,
key maps, gutter markers, bookmarks, and annotation marks.

\section{Concurrency}
The compile lock is acquired before the save operation begins. Rapid repeated
clicks therefore cannot create overlapping save-and-build chains.

\section{Failure recovery}
If saving fails, compilation is cancelled and the conflict remains visible. If
compilation fails, the full log remains available and source diagnostics direct
the user to the first relevant line.

\chapter{Mathematical Material}
Consider a sequence $x_1,\ldots,x_n$. Its cumulative sum is
\begin{equation}
  S_k = \sum_{i=1}^{k}x_i.
\end{equation}
The normalized sequence is
\begin{equation}
  z_i = \frac{x_i-\bar{x}}{s}.
\end{equation}

For repeated measurements indexed by subject $j$ and observation $i$, a simple
random-intercept model can be written as
\begin{equation}
  y_{ij} = \beta_0 + \beta_1 x_{ij} + b_j + \varepsilon_{ij},
  \qquad b_j \sim \mathcal{N}(0,\tau^2).
\end{equation}

\chapter{Soak Log Template}
\section{Session one}
Record the date, browser width, selected theme, split percentage, compilation
result, and any unexpected status transition.

\section{Session two}
Record behavior with two LaTeX tabs, one code tab, and one Markdown tab open at
the same time. Note any shared state that should have remained independent.

\section{Session three}
Record behavior after restarting Atelier. Confirm that the gallery still shows
LaTeX files when the experimental flag is present.

\section{Session four}
Record comments and SyncTeX behavior after multiple recompilations. Pay attention
to stale PDF pages and source highlights that point to the wrong line.

\chapter{Conclusion}
The report provides a durable, editable workload for testing the unified shell.
The shell should become the default only after these workflows remain reliable
across several real sessions.

\appendix
\chapter{Manual Checklist}
\begin{enumerate}
  \item Open both long fixtures from the gallery.
  \item Drag the split separator to several widths.
  \item Toggle the PDF preview off and on.
  \item Compile both documents.
  \item Add and reload anchored comments.
  \item Trigger and clear a compilation error.
  \item Exercise SyncTeX in both directions.
  \item Close one tab while compilation is running.
  \item Restart Atelier and repeat the key actions.
\end{enumerate}

\end{document}
